The ablation results in Table~\ref{tab:ablation} demonstrate the central challenge of evaluating behavioural prompt efficiency. The length-only baseline (using \texttt{Word\_Count} alone) achieves 66.3\% accuracy and a macro-F1 of 0.670, outperforming the full 8-feature text classifier (59.1\% accuracy, 0.576 macro-F1). This confirms that raw prompt length accounts for the largest portion of the variance in the behavioural ground-truth labels, which were constructed using length, conversation tree depth, and repetition penalties.

When length-coupled features (\texttt{unique\_word\_ratio}, \texttt{punct\_density}, and \texttt{avg\_word\_length}) are excluded, accuracy drops to 46.5\% (0.459 macro-F1). Regressing all text features against $\log(\text{Word\_Count})$ and training on residuals yields 66.2\% accuracy (0.667 macro-F1), matching the performance of the length-only baseline. These findings indicate that while stylistic text markers capture domain patterns, prompt length remains the dominant signal under this labelling formulation, and true decoupling from length requires richer syntactic or human-annotated labels.